\documentclass[11pt]{article}
\usepackage[utf8]{inputenc}
\usepackage[T1]{fontenc}
\usepackage[spanish]{babel}
\usepackage{booktabs}
\usepackage{graphicx}
\usepackage{float}
\usepackage{geometry}
\usepackage{amsmath}
\usepackage{array}
\usepackage{longtable}
\geometry{letterpaper, margin=2.5cm}

\title{Implementaci\'on y calibraci\'on secuencial del modelo Black--Litterman}
\author{Proyecto Capstone FinPUC}
\date{Junio 2026}

\begin{document}
\maketitle

\section{Prop\'osito de la calibraci\'on}

La versi\'on anterior del recomendador FinPUC incorporaba Black--Litterman como una extensi\'on natural del modelo media-varianza, pero varios de sus par\'ametros centrales se encontraban fijados de manera heur\'istica: la matriz de \textit{views} $P$, el vector de retornos esperados de las opiniones $Q$, la matriz de incertidumbre $\Omega$ y el escalar $\tau$. Para que la metodolog\'ia fuera defendible en la Entrega 3, se implement\'o una calibraci\'on secuencial basada en evidencia computacional. El objetivo no fue encontrar el \'optimo global de todas las combinaciones posibles, porque esa b\'usqueda ser\'ia costosa y poco trazable, sino construir una heur\'istica ordenada que permitiera explicar cada decisi\'on metodol\'ogica.

La calibraci\'on se ejecut\'o sobre un universo com\'un de 601 activos, con los mismos escenarios de datos utilizados en la entrega anterior: una base de entrenamiento sin pandemia y otra con pandemia. El horizonte principal se mantuvo en $h7\_f2$, es decir, siete a\~nos de historia y dos a\~nos de evaluaci\'on fuera de muestra, con rebalanceo semestral. Esta decisi\'on mantiene la comparabilidad con Markowitz y evita que la calibraci\'on de Black--Litterman se beneficie de una ventana distinta a la usada por el caso base.

\section{Caso base: Markowitz y benchmark}

El punto de partida de la historia metodol\'ogica es el caso base. Markowitz representa la optimizaci\'on financiera pura: estima retornos esperados y covarianzas hist\'oricas, y luego selecciona portafolios seg\'un el perfil de riesgo del cliente. El benchmark equiponderado, en cambio, funciona como control no param\'etrico: no depende de estimaciones de $\mu$ ni de $\Sigma$, por lo que permite medir si la sofisticaci\'on del optimizador realmente agrega valor.

En los resultados previos, Markowitz mostraba un comportamiento ambivalente. En el escenario sin pandemia, el portafolio ``Ideal de mercado'' alcanzaba un Sharpe alto, pero este desempe\~no era sensible a la estimaci\'on hist\'orica. Al incorporar datos de pandemia al entrenamiento, los perfiles m\'as agresivos se deterioraban de forma importante. El benchmark, aunque menos sofisticado, manten\'ia una referencia exigente por su diversificaci\'on simple en activos grandes. Esta tensi\'on justific\'o la introducci\'on de Black--Litterman: no para maximizar retorno bruto a cualquier costo, sino para estabilizar las estimaciones y reducir la dependencia de se\~nales hist\'oricas extremas.

\begin{table}[H]
\centering
\caption{Referencia metodol\'ogica del caso base}
\label{tab:caso_base}
\begin{tabular}{l l r r r}
\toprule
Escenario & Modelo & Retorno total & Drawdown & Sharpe \\
\midrule
Sin pandemia & Markowitz, Ideal de mercado & 76.5\% & -12.6\% & 2.51 \\
Con pandemia & Markowitz, Muy conservador & 54.5\% & -10.2\% & 2.02 \\
Ambos & Benchmark equiponderado & 101.4\% & -21.9\% & 2.05 \\
\bottomrule
\end{tabular}
\end{table}

La lectura de esta tabla no es que un modelo sea universalmente superior. El benchmark logra mayor retorno bruto, pero acepta ca\'idas m\'as profundas. Markowitz puede entregar mejor retorno ajustado por riesgo en ventanas favorables, pero es vulnerable a cambios de r\'egimen. Por ello, la calibraci\'on de Black--Litterman se dise\~n\'o con una regla expl\'icita: una configuraci\'on no pod\'ia ser seleccionada solo por mejorar Sharpe si empeoraba materialmente el drawdown frente a Markowitz.

\section{Camino de calibraci\'on secuencial}

\subsection{Criterio de selecci\'on}

Cada etapa se evalu\'o con un score robusto que combina Sharpe, mejora de Sharpe frente a Markowitz, retorno anual, drawdown, CVaR y estabilidad. Adem\'as, se agreg\'o una regla de resiliencia: si una configuraci\'on empeoraba el drawdown promedio frente a Markowitz en m\'as de tres puntos porcentuales, quedaba degradada en el ranking. Esta regla fue clave, porque una versi\'on inicial de la calibraci\'on eleg\'ia momentum agresivo con alto Sharpe, pero a costa de ca\'idas m\'as severas. El criterio final corrige ese sesgo y alinea el experimento con el objetivo del proyecto: recomendar una metodolog\'ia estable, no una estrategia que solo gane en una m\'etrica.

\subsection{Paso 1: selecci\'on de familia de \textit{view}}

La primera decisi\'on fue determinar qu\'e tipo de informaci\'on deb\'ia entrar al posterior Black--Litterman. Se compararon cuatro familias: momentum general, momentum sobre el top de capitalizaci\'on a seis meses, momentum top/bottom por capitalizaci\'on a un a\~no, y una \textit{view} macro de desempleo. En esta etapa se mantuvieron fijos $\tau = 0.05$, la confianza base de momentum en 0.50 y la confianza de desempleo en 0.35.

El resultado fue metodol\'ogicamente relevante. Momentum general obtuvo el mayor Sharpe medio, pero tambi\'en deterior\'o el drawdown promedio frente a Markowitz. Bajo una lectura puramente oportunista, esa familia habr\'ia sido seleccionada. Sin embargo, al aplicar la regla de resiliencia, la \textit{view} de desempleo se convirti\'o en la alternativa ganadora. Su Sharpe medio fue 1.579 y su mejora promedio frente a Markowitz fue moderada, 0.076 puntos, pero su drawdown promedio fue 2.69 puntos porcentuales mejor que el caso base. La decisi\'on fue, por tanto, privilegiar una se\~nal macro defensiva antes que una se\~nal de momentum con mayor retorno pero peor cola de p\'erdidas.

\begin{table}[H]
\centering
\caption{Selecci\'on de familia de \textit{view}}
\label{tab:familia_view}
\begin{tabular}{l r r r r}
\toprule
Familia & Sharpe medio & $\Delta$ Sharpe & Drawdown medio & Score \\
\midrule
Desempleo macro & 1.579 & 0.076 & -12.9\% & 0.595 \\
Momentum top market-cap 6M & 1.605 & 0.101 & -15.7\% & 0.587 \\
Momentum general & 1.705 & 0.201 & -20.0\% & 0.277 \\
Momentum top/bottom 1Y & 1.489 & -0.014 & -16.4\% & 0.132 \\
\bottomrule
\end{tabular}
\end{table}

\subsection{Paso 2: estructura interna de $P$}

Una vez elegida la familia de desempleo, la calibraci\'on se concentr\'o en la estructura de $P$. En este caso, $P$ no selecciona ganadores y perdedores por ranking de retornos, sino que expresa una opini\'on relativa entre sectores c\'iclicos y defensivos. Con un desempleo asumido menor al neutral, el modelo favorece sectores c\'iclicos frente a sectores defensivos; si el desempleo asumido aumenta sobre el neutral, la direcci\'on de la se\~nal se invierte.

Se evaluaron tres niveles de desempleo asumido, 4\%, 5\% y 6\%, combinados con $\beta$ macro de 0.5, 1.0 y 1.5. La mejor configuraci\'on fue desempleo asumido de 4\% y $\beta=1.5$. Esta elecci\'on no debe interpretarse como una predicci\'on macroecon\'omica exacta, sino como una calibraci\'on de escenario: bajo una econom\'ia con desempleo por debajo de su nivel neutral, una exposici\'on moderada a sectores c\'iclicos mejora el balance retorno-riesgo sin deteriorar la cola de p\'erdidas.

\begin{table}[H]
\centering
\caption{Calibraci\'on de la estructura macro de $P$}
\label{tab:estructura_p}
\begin{tabular}{c c r r r r}
\toprule
Desempleo asumido & $\beta$ macro & Sharpe medio & $\Delta$ Sharpe & Drawdown medio & Score \\
\midrule
4\% & 1.5 & 1.713 & 0.009 & -9.8\% & 0.650 \\
4\% & 1.0 & 1.712 & 0.007 & -9.8\% & 0.606 \\
4\% & 0.5 & 1.710 & 0.006 & -9.7\% & 0.562 \\
5\% & 0.5--1.5 & 1.709 & 0.004 & -9.7\% & 0.510 \\
6\% & 1.5 & 1.704 & -0.000 & -9.7\% & 0.350 \\
\bottomrule
\end{tabular}
\end{table}

\subsection{Paso 3: intensidad de $Q$}

Luego se calibr\'o la intensidad del vector $Q$ mediante un factor de escala $q_{\text{scale}} \in \{0.5,1.0,1.5\}$. Esta etapa responde a una pregunta distinta a la anterior: una vez definida la direcci\'on de la opini\'on macro, cu\'anta fuerza debe tener esa opini\'on en el posterior. El mejor resultado se obtuvo con $q_{\text{scale}}=1.5$, que elev\'o el Sharpe medio a 1.715 y mantuvo un drawdown promedio cercano a -9.8\%.

La interpretaci\'on es que una \textit{view} macro demasiado tenue casi no desplaza el prior de equilibrio, mientras que una intensidad mayor entrega una correcci\'on visible sin generar el deterioro de drawdown observado en momentum. Por eso se seleccion\'o $q_{\text{scale}}=1.5$: no por maximizar retorno aislado, sino porque incrementa la contribuci\'on de la se\~nal manteniendo estable el perfil de p\'erdidas.

\subsection{Paso 4: confianza de $\Omega$}

La matriz $\Omega$ se calibr\'o mediante el nivel de confianza de la \textit{view}. Se probaron valores de 0.20, 0.35 y 0.50. El ganador fue 0.20, es decir, una confianza baja. Este resultado es coherente con la naturaleza de la se\~nal: el desempleo usado es un escenario macro asumido, no una serie externa observada en tiempo real dentro del modelo. Por tanto, no conviene imponerlo con la misma fuerza que una se\~nal de mercado directamente medida.

La baja confianza permite que Black--Litterman conserve el ancla del prior de equilibrio y use la \textit{view} macro como correcci\'on suave. En t\'erminos pr\'acticos, esta etapa evita que el modelo se convierta en una apuesta sectorial excesiva. El Sharpe medio aument\'o a 1.722 y el drawdown medio se mantuvo bajo control, cerca de -9.9\%.

\subsection{Paso 5: ajuste final de $\tau$}

Finalmente se evalu\'o $\tau \in \{0.01,0.025,0.05,0.10,0.20\}$ manteniendo fija la \textit{view} de desempleo, $q_{\text{scale}}=1.5$ y confianza 0.20. El mejor resultado fue $\tau=0.20$, con Sharpe medio 1.730, mejora promedio de 0.025 puntos de Sharpe frente a Markowitz y drawdown promedio de -9.8\%.

La elecci\'on de $\tau$ debe leerse como una sensibilidad final del prior. En esta implementaci\'on, $\Omega$ se calcula en funci\'on de $P\tau\Sigma P^\top$ y la confianza, por lo que $\tau$ no opera como un simple multiplicador de la \textit{view}; tambi\'en afecta la covarianza posterior y, por lo tanto, la estabilidad de la optimizaci\'on. La calibraci\'on muestra que un $\tau$ mayor entrega una versi\'on levemente m\'as flexible del prior sin romper la restricci\'on de resiliencia.

\begin{table}[H]
\centering
\caption{Configuraci\'on final calibrada}
\label{tab:config_final}
\begin{tabular}{l l}
\toprule
Componente & Valor final \\
\midrule
Familia de \textit{view} & Desempleo macro asumido \\
Desempleo asumido & 4\% \\
Desempleo neutral & 5\% \\
$\beta$ macro & 1.5 \\
$q_{\text{scale}}$ & 1.5 \\
Confianza para $\Omega$ & 0.20 \\
$\tau$ & 0.20 \\
Criterio de selecci\'on & Score robusto con penalizaci\'on de drawdown \\
\bottomrule
\end{tabular}
\end{table}

\section{Robustez y escenarios complejos}

El modelo calibrado final no debe presentarse como una metodolog\'ia que gana en todos los escenarios. Su aporte es m\'as espec\'ifico: reduce la severidad de las ca\'idas en escenarios donde Markowitz es sensible al error de estimaci\'on. Esto se observa con claridad en el escenario con pandemia. Para el perfil Arriesgado, Black--Litterman calibrado alcanza Sharpe 1.71 frente a 1.17 de Markowitz, una mejora de 46.7\%, y reduce el drawdown desde -19.1\% a -11.9\%. Para el perfil Muy arriesgado, el Sharpe sube de 0.71 a 1.08 y el drawdown mejora desde -29.2\% a -24.7\%. En el perfil Neutro, el Sharpe aumenta de 1.66 a 1.79 y el drawdown mejora desde -11.1\% a -9.3\%.

\begin{table}[H]
\centering
\caption{Modelo calibrado frente a Markowitz en escenario con pandemia}
\label{tab:pandemia}
\begin{tabular}{l r r r r r}
\toprule
Perfil & Sharpe BL & Sharpe Markowitz & Mejora & Drawdown BL & Drawdown Markowitz \\
\midrule
Neutro & 1.79 & 1.66 & 7.8\% & -9.3\% & -11.1\% \\
Arriesgado & 1.71 & 1.17 & 46.7\% & -11.9\% & -19.1\% \\
Muy arriesgado & 1.08 & 0.71 & 51.7\% & -24.7\% & -29.2\% \\
Conservador & 1.78 & 1.78 & 0.2\% & -9.0\% & -8.9\% \\
Muy conservador & 1.77 & 1.79 & -0.9\% & -8.9\% & -8.8\% \\
\bottomrule
\end{tabular}
\end{table}

La validaci\'on robusta confirma la misma idea con un matiz importante. En ventanas m\'oviles para el perfil Neutro, Black--Litterman calibrado obtuvo Sharpe medio 1.44, inferior al 1.63 de Markowitz, pero su drawdown medio fue -11.1\%, mejor que el -13.5\% del caso base. Por lo tanto, el modelo calibrado no es superior si se eval\'ua \'unicamente por Sharpe promedio en todas las ventanas. Su fortaleza est\'a en la contenci\'on de p\'erdidas: sacrifica parte del retorno esperado para reducir la exposici\'on a ca\'idas m\'as profundas.

Esta distinci\'on es central para la conclusi\'on del proyecto. Si el objetivo fuera maximizar Sharpe promedio en cualquier condici\'on, el modelo no ser\'ia una mejora universal. Si el objetivo es construir un recomendador m\'as resiliente para clientes con sensibilidad a p\'erdidas, la configuraci\'on calibrada es defendible, especialmente en perfiles donde Markowitz se deteriora ante informaci\'on de estr\'es.

\section{Integraci\'on posterior con Monte Carlo P4}

Una vez calibrado el portafolio, sus KPIs fueron usados como entrada para la simulaci\'on Monte Carlo P4. Esta integraci\'on se realiz\'o despu\'es de la optimizaci\'on financiera. La utilidad de la empresa no entr\'o en la funci\'on objetivo del optimizador de portafolios; se calcul\'o posteriormente como consecuencia del saldo administrado, la aceptaci\'on de recomendaciones y el retiro del cliente. Esta separaci\'on es metodol\'ogicamente necesaria: el portafolio se optimiza para el cliente, mientras que la empresa eval\'ua despu\'es la sostenibilidad econ\'omica de esa recomendaci\'on.

Los resultados P4 muestran que la configuraci\'on calibrada no domina comercialmente a Markowitz en todos los casos. El mejor score P4 observado corresponde a Markowitz en el perfil Neutro sin pandemia, con riqueza media USD 2,813, retiro 0.9\%, utilidad empresa USD 194 y score 2,998. En comparaci\'on, Black--Litterman calibrado obtiene en el perfil Neutro sin pandemia riqueza media USD 2,306, retiro 0.9\%, utilidad empresa USD 145 y score 2,442. La diferencia se explica porque el modelo calibrado es m\'as defensivo: reduce drawdown, pero tambi\'en disminuye retorno esperado y, por ende, saldo administrado futuro.

\begin{table}[H]
\centering
\caption{Evaluaci\'on posterior P4: cliente y empresa}
\label{tab:p4}
\begin{tabular}{l l l r r r r}
\toprule
Modelo & Escenario & Perfil & Riqueza & Retiro & Utilidad & Score P4 \\
\midrule
Markowitz & Sin pandemia & Neutro & 2,813 & 0.9\% & 194 & 2,998 \\
Markowitz & Sin pandemia & Conservador & 2,533 & 14.4\% & 198 & 2,587 \\
Markowitz & Con pandemia & Neutro & 2,389 & 0.7\% & 154 & 2,536 \\
BL calibrado & Sin pandemia & Conservador & 2,457 & 14.9\% & 194 & 2,501 \\
BL calibrado & Con pandemia & Conservador & 2,393 & 13.7\% & 189 & 2,445 \\
BL calibrado & Con pandemia & Neutro & 2,278 & 0.4\% & 142 & 2,417 \\
\bottomrule
\end{tabular}
\end{table}

La simulaci\'on P4, por tanto, no contradice la calibraci\'on. M\'as bien, aclara el uso del modelo. Black--Litterman calibrado no debe venderse como el m\'aximo generador de riqueza terminal. Su funci\'on es distinta: entregar una alternativa m\'as conservadora en t\'erminos de drawdown para perfiles donde la permanencia del cliente y la tolerancia a p\'erdidas son variables relevantes. La utilidad de FinPUC se mantiene como una m\'etrica posterior, y su menor valor relativo en algunos casos refleja precisamente el costo de una recomendaci\'on m\'as defensiva.

\section{Conclusi\'on}

La calibraci\'on secuencial permiti\'o transformar Black--Litterman desde un conjunto de par\'ametros elegidos heur\'isticamente a una metodolog\'ia defendible. El camino de calibraci\'on mostr\'o que una selecci\'on basada solo en Sharpe habr\'ia favorecido momentum general; sin embargo, esa alternativa empeoraba el drawdown promedio frente al caso base. Al incorporar expl\'icitamente una penalizaci\'on por deterioro de p\'erdidas, el modelo final recomendado pas\'o a ser Black--Litterman con \textit{view} macro de desempleo, desempleo asumido de 4\%, desempleo neutral de 5\%, $\beta=1.5$, $q_{\text{scale}}=1.5$, confianza 0.20 y $\tau=0.20$.

La recomendaci\'on principal es utilizar esta configuraci\'on calibrada como metodolog\'ia defensiva para perfiles Neutro, Arriesgado y Muy arriesgado cuando el entrenamiento incorpora escenarios complejos como la pandemia. En el escenario con pandemia, el modelo mejora el Sharpe del perfil Neutro en 7.8\%, del perfil Arriesgado en 46.7\% y del perfil Muy arriesgado en 51.7\% frente a Markowitz. Adem\'as, reduce el drawdown en esos tres perfiles: de -11.1\% a -9.3\% en Neutro, de -19.1\% a -11.9\% en Arriesgado, y de -29.2\% a -24.7\% en Muy arriesgado. Esta evidencia respalda una conclusi\'on fuerte pero acotada: Black--Litterman calibrado no reemplaza a Markowitz en todos los contextos, pero s\'i entrega una mejora robusta para perfiles expuestos al deterioro de estimaciones hist\'oricas bajo estr\'es.

Para perfiles conservadores, la recomendaci\'on debe ser m\'as cuidadosa. En Muy conservador y Conservador las diferencias de Sharpe frente a Markowitz son peque\~nas, y en algunos casos no compensan la complejidad adicional del modelo. En t\'erminos P4, Markowitz mantiene mejores scores en varios escenarios porque genera mayor riqueza terminal y, por lo tanto, mayor saldo administrado. Por ello, la conclusi\'on de negocio no es que Black--Litterman calibrado maximice siempre la utilidad de FinPUC, sino que ofrece una alternativa metodol\'ogicamente m\'as resiliente cuando el objetivo prioritario es controlar p\'erdidas y sostener una recomendaci\'on defendible ante escenarios adversos.

En s\'intesis, la Entrega 3 debe presentar esta calibraci\'on como una historia de aprendizaje: Markowitz establece el caso base, momentum revela el riesgo de perseguir Sharpe sin control de drawdown, la \textit{view} de desempleo introduce una se\~nal macro m\'as estable, y la calibraci\'on de $P$, $Q$, $\Omega$ y $\tau$ permite llegar a una configuraci\'on final trazable. La recomendaci\'on final no es universal, sino segmentada: Black--Litterman calibrado con desempleo es recomendable para perfiles de mayor exposici\'on al riesgo en escenarios de estr\'es, mientras que Markowitz y el benchmark siguen siendo referencias necesarias para evaluar riqueza terminal y desempe\~no comercial.

\end{document}
